\section*{अध्याय \ref{ch:b1:filters-transfer-functions} --- फ़िल्टर और स्थानांतरण फलन}

\begin{solution}{exo:b1:filters-transfer-functions:1}
$0.5 \to \qty{-6.0}{dB}$; $10 \to \qty{20}{dB}$; $1/\sqrt2 \to \qty{-3.0}{dB}$;
$100 \to \qty{40}{dB}$। $\qty{-40}{dB} \to 0.01$; $\qty{+6}{dB} \to 2.0$।
शृंखलन: $-6 - 20 = \qty{-26}{dB}$, यानी अनुपात $0.05$।
\end{solution}

\begin{solution}{exo:b1:filters-transfer-functions:2}
$f_c = 1/(2\pi RC) = 1/(2\pi \times 10^4 \times 1.6\times10^{-8}) =
\qty{1.0}{kHz}$। लब्धियाँ $-10\log(1 + (f/f_c)^2)$: \qty{100}{Hz}:
\qty{-0.04}{dB}; \qty{1}{kHz}: \qty{-3}{dB}; \qty{10}{kHz}: \qty{-20}{dB};
\qty{100}{kHz}: \qty{-40}{dB}।
\end{solution}

\begin{solution}{exo:b1:filters-transfer-functions:3}
(क) निम्न-पारक ($C$ ऊँची $f$ पर निर्गम को लघुपथित कर देता है); (ख) उच्च-पारक
($C$ कम $f$ पर खुल जाता है); (ग) द्वितीय कोटि का निम्न-पारक ($L$ ऊँची $f$ पर
खुलता है, $C$ लघुपथित करता है); (घ) बैंड-पारक ($C$ कम $f$ को रोकता है, $L$ ऊँची $f$ को)।
\end{solution}

\begin{solution}{exo:b1:filters-transfer-functions:4}
$4E/n\pi$: \qty{1}{kHz} पर \qty{6.37}{V}, \qty{3}{kHz} पर \qty{2.12}{V},
\qty{5}{kHz} पर \qty{1.27}{V}। वर्ग तरंग आधे आवर्तकाल के खिसकाव के
सापेक्ष प्रतिसममित है ($s(t + T/2) = -s(t)$), और सम संनादी ऐसे
नहीं होते।
\end{solution}

\begin{solution}{exo:b1:filters-transfer-functions:5}
\qty{100}{Hz} पर और उसके उपयोगी संनादियों पर $\omega \ll \omega_c$:
अतः $\underline H \approx j\omega/\omega_c$, यानी अवकलज बटा
$\omega_c$। त्रिभुज की ढाल $\pm 4Ef = \pm\qty{400}{V/s}$ है, अतः
निर्गम \omterm{def:b1:wave-propagation:signal}{आयाम} $400/(2\pi \times 10^4) =
\qty{6.4}{mV}$ की वर्ग तरंग है।
\end{solution}

\begin{solution}{exo:b1:filters-transfer-functions:6}
$\omega \gg \omega_c$: $u_s \approx (1/\tau)\int u_e\,\dd t$, जहाँ $\tau =
1/(2\pi \times 10) = \qty{16}{ms}$। $\pm E$ का समाकल ऐसी
त्रिभुज है जो प्रति अर्ध-आवर्तकाल $ET/2$ चढ़ती है: अतः झूला $ET/2\tau$, और \omterm{def:b1:wave-propagation:signal}{आयाम}
$ET/4\tau = 1.0 \times 10^{-3}/(4 \times 0.016) = \qty{16}{mV}$ शून्य के
आसपास।
\end{solution}

\begin{solution}{exo:b1:filters-transfer-functions:7}
$f_0 = \qty{5.03}{kHz}$। $Q = \sqrt{L/C}/R = 1/\sqrt2$: अतः $R = \sqrt2 \times
316 = \qty{447}{\ohm}$। $f_0$ पर: \qty{-3}{dB}; और $10f_0$ पर: $1/\sqrt{1 +
10^4} \to \qty{-40}{dB}$।
\end{solution}

\begin{solution}{exo:b1:filters-transfer-functions:8}
$G = 1/\sqrt{1 + Q^2(x - 1/x)^2}$, जहाँ $x = f/3$: \qty{1}{kHz}: $x = 1/3$,
$Q(x - 1/x) = -53$, $G = 0.019$; \qty{3}{kHz}: $G = 1$; \qty{5}{kHz}:
$x = 5/3$, $21$, $G = 0.047$। निर्गम: \qty{3}{kHz} पर \qty{2.1}{V}, साथ में
\qty{1}{kHz} पर \qty{0.12}{V} और \qty{5}{kHz} पर \qty{0.06}{V} --- यानी
लगभग शुद्ध \qty{3}{kHz} की ज्या।
\end{solution}

\begin{solution}{exo:b1:filters-transfer-functions:9}
संधियाँ: $a = jRC\omega$ के साथ, पहला संधारित्र अनुप्रवाह से पहले $R$ और
अनुप्रवाह में $R + 1/jC\omega$ देखता है: अतः $\underline H = 1/(1 + 3a + a^2)$, यानी
$1/[1 - (\omega/\omega_c)^2 + 3j\omega/\omega_c]$। $\omega_c$ पर: $1/3j$,
$G = 1/3$, \qty{-9.5}{dB}; जबकि अलग-अलग दोनों लब्धियों का गुणनफल
$1/2$, यानी \qty{-6}{dB} होता। $f_c = \qty{1.0}{kHz}$।
\end{solution}

\begin{solution}{exo:b1:filters-transfer-functions:10}
$\omega_1$ से नीचे: \qty{0}{dB}; $\omega_1$ और $\omega_2$ के बीच: प्रति दशक
\qty{20}{dB} चढ़ती हुई; $\omega_2$ से ऊपर: $20\log(\omega_2/\omega_1) =
\qty{40}{dB}$ पर समतल। \omterm{def:b1:wave-propagation:signal}{कला} $\arctan(\omega/\omega_1) - \arctan(\omega/\omega_2)$:
$\omega = 10^3$ पर $\arctan 10 - \arctan 0.1 = 84^\circ - 6^\circ =
78^\circ$।
\end{solution}

\begin{solution}{exo:b1:filters-transfer-functions:11}
कटऑफ़ से कहीं ऊपर $G \approx f_c/f$; अतः फ़िल्टरन के बाद तरंगिका का अनुपात
है
\[
\frac{(4E/3\pi)(f_c/100)}{2E/\pi} = \frac{2}{3}\,\frac{f_c}{100} < 0.01,
\]
अतः $f_c < \qty{1.5}{Hz}$ और $\tau = 1/2\pi f_c > \qty{0.1}{s}$। भार स्वयं ही
$R$ है: उसके आरपार लगा संधारित्र न विभव-पात की क़ीमत लेता है न शक्ति की,
जबकि श्रेणी में लगा कोई \omterm{def:b1:dc-circuits:resistor}{प्रतिरोधक} दोनों बरबाद करता।
\end{solution}

\begin{solution}{exo:b1:filters-transfer-functions:12}
$f_c = 1/(2\pi RC) = 1/(2\pi \times 1.6\times10^5 \times 10^{-7}) =
\qty{10}{Hz}$, $\tau = RC = \qty{16}{ms}$। किसी समतल चोटी के दौरान
संधारित्र $R$ से होकर आवेशित होता है: अतः निर्गम हर किनारे से
$\eu^{-t/\tau}$ की तरह क्षीण होता है। $T/2$ पर झोल: $1 - \eu^{-T/2\tau}$: \qty{50}{Hz}
($T/2 = \qty{10}{ms}$): $1 - \eu^{-0.63} = 47\%$; \qty{1}{kHz}
($T/2 = \qty{0.5}{ms}$): $3\%$।
\end{solution}

\begin{solution}{pb:b1:filters-transfer-functions:1}
\textbf{1.} $\underline H = \underline U_{\text{चालक}}/\underline U_e$;
$G_{\mathrm{dB}} = 20\log\abs{\underline H}$।

\textbf{2.} $\underline H_L = R/(R + jL\omega) = 1/(1 + j\omega/\omega_c)$,
$\omega_c = R/L$.

\textbf{3.} $L = R/\omega_c = 8.0/(2\pi \times 2000) = \qty{0.64}{mH}$.

\textbf{4.} $\underline H_H = R/(R + 1/jC\omega) = jRC\omega/(1 + jRC\omega)$,
$\omega_c = 1/RC$; $C = 1/(R\omega_c) = \qty{9.9}{\micro F} \approx
\qty{10}{\micro F}$।

\textbf{5.} $f_c$ पर: हर एक को \qty{-3}{dB}। \qty{200}{Hz} ($x = 0.1$): वूफ़र
\qty{-0.04}{dB}, ट्वीटर \qty{-20}{dB}। \qty{20}{kHz} ($x = 10$): वूफ़र
\qty{-20}{dB}, ट्वीटर \qty{-0.04}{dB}।

\textbf{6.} $G_L^2 + G_H^2 = (1 + x^2)/(1 + x^2) = 1$: अतः बराबर भारों के साथ
दोनों शक्तियाँ जुड़कर उतनी ही बनती हैं जितनी अकेला \qty{8}{\ohm} का चालक
लेता --- यानी बँटी हुई, खोई हुई नहीं।

\textbf{7.} $-45^\circ$ (वूफ़र) और $+45^\circ$ (ट्वीटर): यानी $90^\circ$
का अंतर; और $f_c$ पर दोनों शंकु साथ-साथ नहीं चलते।

\textbf{8.} प्रति अष्टक \qty{\pm 6}{dB} (यानी प्रति दशक \qty{20}{dB})।

\textbf{9.} $jL\omega$ और $R \parallel 1/jC\omega =
R/(1 + jRC\omega)$ के बीच विभाजक: $\underline H = R/[R + jL\omega(1 + jRC\omega)] =
1/(1 - LC\omega^2 + jL\omega/R)$।

\textbf{10.} $\omega_0 = 1/\sqrt{LC}$; और $L\omega/R = x/Q$ से $Q =
R/L\omega_0 = R\sqrt{C/L}$।

\textbf{11.} $\abs{\underline H}^{-2} = (1 - x^2)^2 + x^2/Q^2 = 1 + x^4 +
x^2(1/Q^2 - 2) = 1 + x^4$ जब $Q^2 = \tfrac12$; और $x = 1$ पर $G = 1/\sqrt2$।

\textbf{12.} $L = R/(Q\omega_0) = 8\sqrt2/(1.257\times10^4) = \qty{0.90}{mH}$;
$C = 1/(L\omega_0^2) = \qty{7.0}{\micro F}$.

\textbf{13.} $x = 2$: $1/\sqrt{17}$, \qty{-12.3}{dB}; $x = 4$: $1/\sqrt{257}$,
\qty{-24.1}{dB}: यानी प्रति अष्टक \qty{12}{dB}।

\textbf{14.} $L$ और $C$ की अदला-बदली $x \leftrightarrow 1/x$ बदल देती है (और
$Q$ वही रहता है): अतः $\underline H_H = -x^2/(1 - x^2 + jx/Q)$।

\textbf{15.} $x = 1$ पर: $\underline H_L = Q/j = -jQ$ ($-90^\circ$),
$\underline H_H = jQ$ ($+90^\circ$): अतः क्रॉसओवर पर दोनों चालक \omterm{def:b1:wave-propagation:signal}{कला} में
विपरीत हैं और ध्वनिकीय रूप से एक-दूसरे को काट देते; इसलिए ट्वीटर की तारें
उलटी ध्रुवता में जोड़ी जाती हैं।

\textbf{16.} \qty{2}{kHz} से ऊपर के संनादी: $n \geq 10$ (\qty{2.2}{kHz}
और उससे ऊपर) अधिकतर ट्वीटर को जाते हैं।

\textbf{17.} $9 \times 220 = \qty{1980}{Hz} \approx f_c$: अतः बराबर बँटता है,
हर एक को \qty{-3}{dB}।

\textbf{18.} वूफ़र ($G_L$): $x = 0.5$: $0.89$; $x = 1.5$: $0.55$; $x = 2.5$:
$0.37$। ट्वीटर ($G_H$): $0.45$, $0.83$, $0.93$।

\textbf{19.} ट्वीटर: उसका श्रेणी संधारित्र दिष्ट धारा रोक देता है। वूफ़र
विस्थापन झेलता है (\qty{8}{\ohm} के आरपार \qty{1}{V}: यानी \qty{125}{mA},
\qty{0.125}{W} की ऊष्मा और खिसका हुआ शंकु)।

\textbf{20.} प्रथम कोटि, $x = 0.25$: $G_H = 0.25/\sqrt{1.0625} = 0.24$
(\qty{-12}{dB}), शक्ति-अंश $6\%$। द्वितीय कोटि: $0.0625/\sqrt{1 +
0.0039} = 0.06$ (\qty{-24}{dB}), $0.4\%$। कुछ ही वाट का ट्वीटर प्रथम कोटि
वाली अभिकल्पना में किसी ऊँचे मंद्र स्वर का $6\%$ पाता है --- और वह
जल सकता है।

\textbf{21.} $\underline Z = 8 + jL_{vc}\omega$: \qty{2}{kHz}: $8 + 6.3j$,
$\abs{\underline Z} = \qty{10.2}{\ohm}$; \qty{8}{kHz}: $8 + 25.1j$,
\qty{26.3}{\ohm}।

\textbf{22.} $\underline H = \underline Z/(\underline Z + jL\omega)$:
\qty{2}{kHz}: $10.2/\abs{8 + 14.3j} = 10.2/16.4 = 0.62$ (\qty{-4.2}{dB});
\qty{8}{kHz}: $26.3/\abs{8 + 57.3j} = 0.45$ (\qty{-6.9}{dB}), न कि
\qty{-12.3}{dB}: यानी कुंडली की बढ़ती \omterm{def:b1:sinusoidal-impedance:impedance}{प्रतिबाधा} श्रेणी वाली कुंडली से लड़ती है,
और ढाल घटकर प्रति अष्टक कुछ ही \omterm{def:b1:filters-transfer-functions:bode}{डेसिबल} रह जाती है।

\textbf{23.} प्रवेश्यताएँ: $1/(R + jL_{vc}\omega) + jC_z\omega/(1 + jRC_z\omega)$;
और $C_z = L_{vc}/R^2$ के साथ दूसरा पद $jL_{vc}\omega/[R(R +
jL_{vc}\omega)]$ है, अतः योग $(R + jL_{vc}\omega)/[R(R + jL_{vc}\omega)]
= 1/R$। और $C_z = 0.50\times10^{-3}/64 = \qty{7.8}{\micro F}$।

\textbf{24.} $U = \sqrt{PR} = \sqrt{400} = \qty{20}{V}$ प्रभावी। $f_c$ पर हर
चालक को आधी: \qty{25}{W}। और \qty{500}{Hz} पर ट्वीटर को $6\%$:
\qty{3}{W}।

\textbf{25.} प्रथम कोटि: दो अवयव, शक्ति में पूरक, प्रति अष्टक कोमल
\qty{6}{dB} ढालें, चालकों के बीच $90^\circ$, और मंद्र स्वर का ट्वीटर में
रिसाव। द्वितीय कोटि: चार अवयव, प्रति अष्टक \qty{12}{dB}, क्रॉसओवर पर
$180^\circ$ (ट्वीटर को उलटिए), और ट्वीटर की कहीं बेहतर सुरक्षा। और दोनों ही
हालत में अभिकल्पना तभी टिकती है जब वूफ़र के प्रेरकत्व को किसी ज़ोबेल जाल से
साध लिया जाए।
\end{solution}
